\chapter{Perpetual futures and funding}\label{ch:perp}

\noindent\emph{Nothing here is investment advice. Venue rules change without notice; anything quoted
as a rule is dated and re-checkable.}\medskip

\noindent The perpetual future---``perp''---is the dominant crypto derivative: a futures-style
contract with \emph{no expiry date}. It is where most leveraged price exposure lives, where most
liquidations happen, and where the funding rate acts as a continuous toll on crowded positions.
This chapter is about the mechanics: what a perp is, how funding replaces expiry convergence,
how margin and liquidation actually compute, and how to read funding as data.

\section{What a perp is and is not}\label{sec:whatperp}

An expiring future needs no tether to spot: at expiry the contract \emph{settles} to the spot
price, so arbitrage forces convergence as the date approaches. A perp never settles, so nothing
pulls its price toward the index except a manufactured cost---\emph{funding}.

\begin{principle}[The two engines]
A perpetual future is a leveraged position on an index price with two moving parts:
\textbf{(i)} a funding transfer between longs and shorts that keeps the perp near the index,
and \textbf{(ii)} a margin/liquidation engine that closes positions whose collateral can no
longer cover their losses. Everything else---leverage sliders, mark price, insurance funds---is
in service of these two.
\end{principle}

\begin{figure}[htbp]\centering
\begin{tikzpicture}[x=0.62cm,y=0.055cm]
  \draw[->,slate] (0,0) -- (12.6,0) node[below left=-1pt,font=\scriptsize]{days};
  \draw[->,slate] (0,-30) -- (0,32) node[above left=-1pt,font=\scriptsize]{$F-S$ (bp)};
  \draw[slate,dashed] (0,0)--(12.4,0);
  % expiring futures: decay to zero
  \foreach \x/\b in {0/26,2/19,4/12,6/7,8/3.5,10/1}
    \draw[accent,very thick] (\x,\b) -- ({\x+2},{\b*0.55});
  % perp: oscillates, never expires
  \draw[amber,very thick] plot coordinates {(0,22)(1.5,14)(3,19)(4.5,8)(6,15)(7.5,6)(9,13)(10.5,4)(12,9)};
  \node[anchor=west,font=\scriptsize,text=accent] at (12.4,6) {expiring futures};
  \node[anchor=west,font=\scriptsize,text=amber] at (12.4,12) {perp};
  \foreach \x in {2,4,6,8,10,12} \draw[slate] (\x,0.6)--(\x,-0.6) node[below,font=\scriptsize]{\x};
  \node[align=center,font=\scriptsize,text=slate] at (6,-22)
    {expiry forces convergence to spot; the perp has no expiry,\\ so funding must be manufactured to pull it back};
\end{tikzpicture}

\caption{Convergence. Expiring futures are pulled to spot at expiry; a perp oscillates around
the index forever because funding, not expiry, does the pulling.}
\label{fig:perpexpiry}
\end{figure}

Concrete numbers make the size of the instrument clear. On a \$50{,}000 index, a trader going
long 0.5\,BTC of perp at $25\times$ leverage posts about \$1{,}000 of margin. The position is
\$25{,}000 of notional. A $1\%$ index move is $\pm\$250$ of PnL---$25\%$ of posted margin.
That asymmetry is the whole appeal and the whole danger.

\section{The funding mechanism}\label{sec:funding}

Every venue implements the same skeleton: measure the perp's \emph{premium} over the index,
turn it into a rate $f_t$, and transfer $f_t$ per unit of notional between longs and shorts each
funding epoch. Positive $f_t$: longs pay shorts. Negative: shorts pay longs.

\begin{figure}[htbp]\centering
\begin{tikzpicture}[every node/.style={draw=slate,fill=white,rounded corners=1pt,
  font=\scriptsize,align=center,inner sep=4pt}]
  \node (prem) {premium $p_t$\\ perp vs index};
  \node (clamp) [right=1.4cm of prem] {clamp\\ $|f_t|\le f_{\max}$};
  \node (split) [right=1.4cm of clamp] {sign of $f_t$};
  \node (long) [above right=0.5cm and 1.2cm of split,text=good] {longs pay shorts\\ $f_t>0$};
  \node (short) [below right=0.5cm and 1.2cm of split,text=amber] {shorts pay longs\\ $f_t<0$};
  \draw[->,ink] (prem) -- (clamp);
  \draw[->,ink] (clamp) -- (split);
  \draw[->,ink] (split) -- (long);
  \draw[->,ink] (split) -- (short);
\end{tikzpicture}

\caption{The funding pipeline: premium $\to$ rate (clamped) $\to$ peer-to-peer transfer.}
\label{fig:fundflow}
\end{figure}

\begin{principle}[funding rate, schematic]
Let $p_t$ be the perp's premium index (e.g.\ impact-mid minus index, smoothed). The funding rate
is typically $f_t = \mathrm{clamp}(\tilde p_t,\,-f_{\max},\,f_{\max}) + I$ where $I$ is a fixed
interest component. On Hyperliquid, funding is paid \textbf{every hour} and the interest term is
about $0.125$\,bp/h ($=0.01\%$ per $8$h) \emph{[unverified: venue rules, checked 2026-10-11,
re-check quarterly]}.
\end{principle}

We fetched 500 hours of real BTC funding from the Hyperliquid public API
(\texttt{fundingHistory}, no auth) and pinned the response to \texttt{data/hl\_funding\_btc.json}:

\begin{lstlisting}[language=Python, numbers=none]
d=json.load(open("data/hl_funding_btc.json"))
f=np.array([float(x["fundingRate"]) for x in d])
\end{lstlisting}
\begin{lstlisting}[numbers=none]
500 hourly funding prints  1757548800090 .. 1759345200083
funding rate: mean +0.113bp/h  sd 0.038bp  min -0.16bp  max +0.12bp
f>0: 97% of hours   f=+0.125bp exact: 85%
premium: mean -1.0bp   mean(f - premium) = +1.08bp ~ the 1.25bp/8h interest floor
annualised carry at mean rate: +9.9%
30d-equivalent carry: 56bp of notional over 500h
\end{lstlisting}

Read the structure out of the data: $85\%$ of hours print exactly the $0.125$bp/h interest
component, meaning the clamped premium term was near zero---a quiet, balanced book. Yet the
mean rate $+$1.08bp over the premium says the interest floor, not the premium, set the price of
being long for that window.

\begin{figure}[htbp]\centering
\begin{tikzpicture}
\begin{axis}[width=0.86\linewidth,height=0.32\linewidth,
    xlabel={hour},ylabel={funding (bp/h)},
    xmin=0,xmax=500,ymin=-0.2,ymax=0.2,
    tick style={draw=none},axis line style={draw=slate},
    xtick={0,100,200,300,400,500},font=\scriptsize]
  \addplot[accent,thick] file{figures/data/F_FUND.tex};
  \addplot[slate,dashed] coordinates {(0,0.125)(500,0.125)};
  \addplot[slate,dotted] coordinates {(0,0)(500,0)};
  \node[font=\scriptsize,text=slate] at (axis cs:250,0.15) {$+0.125$bp interest floor};
\end{axis}
\end{tikzpicture}

\caption{Real BTC funding on Hyperliquid, 500 hours ending 2026-10-09. The dashed line is the
$0.125$bp/h interest component. Data: \texttt{data/hl\_funding\_btc.json}.}
\label{fig:fundts}
\end{figure}

\begin{example}\label{ex:fundpay}
\textbf{Funding payment.} You are long $2$\,BTC notional at index \$50{,}000. The next three
hourly funding rates print $+0.5$\,bp, $+0.12$\,bp, $-0.3$\,bp. What do you pay or receive,
and what does a sustained $+0.5$\,bp/h regime cost per day?
\begin{solution}
Payment each epoch $=$ position $\times$ notional $\times f_t$:
$2\times 50{,}000 \times (0.5+0.12-0.3)\times 10^{-4}
 = 100{,}000\times 3.2\times 10^{-5} = \$3.2$ paid over the three hours.
At a sustained $+0.5$\,bp/h the day cost is
$100{,}000\times 0.5\times 10^{-4}\times 24 = \$120$/day---about $8.8\%$ per month on the
\$10{,}000 margin you might have posted at $10\times$ leverage.
\end{solution}
\end{example}

\section{Basis and carry}\label{sec:basis}

The \emph{basis} $B_t = F_t - S_t$ (or $F_t/S_t - 1$) is the perp's premium over spot. For an
expiring contract basis decays to zero on a schedule; for a perp it is whatever funding tolerates.
The annualised basis $\hat b = B/S \times (365/D)$ compares carry opportunities across tenors.

\begin{principle}[cash-and-carry]
Long spot $+$ short perp in equal notional is \emph{delta-neutral}: price moves cancel. The
position's return stream is funding: shorts collect $f_t$ when $f_t>0$. Its risks are
liquidation on the short leg, funding flipping negative, and basis widening \emph{before}
collapsing---you can be right about the level and still be forced out.
\end{principle}

\begin{figure}[htbp]\centering
\begin{tikzpicture}
\begin{axis}[width=0.86\linewidth,height=0.32\linewidth,
    xlabel={hour},ylabel={perp basis (bp)},
    xmin=0,xmax=720,ymin=-60,ymax=60,
    tick style={draw=none},axis line style={draw=slate},
    xtick={0,180,360,540,720},font=\scriptsize]
  \addplot[accent,thick] file{figures/data/B_BASIS.tex};
  \addplot[slate,dotted] coordinates {(0,0)(720,0)};
  \addplot[slate,dashed] coordinates {(0,20)(720,20)};
  \addplot[slate,dashed] coordinates {(0,-20)(720,-20)};
  \node[font=\scriptsize,text=slate] at (axis cs:60,32) {$\pm20$bp};
\end{axis}
\end{tikzpicture}

\caption{Simulated mean-reverting basis on a GBM index (seed 11): sd $\approx19$\,bp, inside
$\pm20$\,bp $70\%$ of hours. \texttt{basis\_carry.py}.}
\label{fig:basists}
\end{figure}

\begin{example}\label{ex:carry}
\textbf{Carry yield.} \texttt{funding\_sim.py} simulates a premium process mean-reverting to
$+0.13$\,bp/h with noise, clamps the rate, and applies it to a 2\,BTC position:
\begin{lstlisting}[numbers=none]
T=720h  mean funding 0.091bp/h  sum 65.78bp  ~APR 800bp
long pays +275 USD over 30d (positive = paid by longs)
price PnL -25,402  funding -275  net -25,677
short-perp carry yield 30d: 27.5bp of notional
\end{lstlisting}
Two lessons. (i) The long's net PnL is price PnL \emph{minus} a funding bill that accrues even
when price is flat---the toll of holding the crowded side. (ii) The flip side is a real trade:
short perp $+$ long spot harvested $27.5$\,bp of notional in a simulated month at these rates.
\begin{solution}
The carry figure is $\sum_t f_t$ times notional, and only the funding flips sign between the
two sides---the delta terms cancel exactly. The printed numbers are not an advertisement: the
same file shows the risk, because the simulation's price drop cost the naked long \$25{,}402.
Carry trades fail in both directions at once: funding flips \emph{and} the short leg's margin
is stressed precisely when basis widens.
\end{solution}
\end{example}

\section{Margin, mark price and liquidation}\label{sec:liq}

Initial margin (IM) is collateral posted to open; maintenance margin (MM) is the floor below
which the venue takes over the position. For an isolated long of quantity $q$ at entry $E$
with margin $w$, liquidation hits when $q(E-P) = w - qP\,\mathrm{MM}$, giving
\begin{equation}\label{eq:liq}
P_{\mathrm{liq}} \;=\; E\,\frac{1-\mathrm{IM}}{1-\mathrm{MM}},
\qquad \mathrm{IM}=\tfrac{w}{qE}.
\end{equation}
At $25\times$ leverage ($\mathrm{IM}=4\%$, $\mathrm{MM}=0.5\%$), $P_{\mathrm{liq}}=E\times
0.96/0.995 \approx 0.9648E$---a $3.5\%$ adverse move, not the naive $1/25=4\%$. The $0.5\%$
maintenance margin is why liquidation always sits closer than ``one over leverage''.

\begin{figure}[htbp]\centering
\begin{tikzpicture}
\begin{axis}[width=0.7\linewidth,height=0.36\linewidth,
    xlabel={leverage ($\times$)},ylabel={drop to liquidation (\%)},
    xmin=0,xmax=52,ymin=0,ymax=55,
    tick style={draw=none},axis line style={draw=slate},
    xtick={1,10,20,30,40,50},font=\scriptsize]
  \addplot[accent,very thick,mark=*] file{figures/data/L_LEV.tex};
  \addplot[slate,dashed] coordinates {(1,0.5)(50,0.5)};
  \node[font=\scriptsize,text=slate,align=left] at (axis cs:33,30)
    {$d_\mathrm{liq} = 1-\frac{1-\mathrm{IM}}{1-\mathrm{MM}}$\\ MM $=0.5\%$};
\end{axis}
\end{tikzpicture}

\caption{Distance to liquidation vs leverage at $\mathrm{MM}=0.5\%$. Above $25\times$, ordinary
hourly volatility reaches your liquidation level within hours.}
\label{fig:liqcurve}
\end{figure}

\begin{example}\label{ex:liq}
\textbf{Liquidation price.} \texttt{liq\_calc.py} computes \eqref{eq:liq} for entry 50{,}000,
0.5\,BTC, $25\times$:
\begin{lstlisting}[numbers=none]
entry 50000 qty 0.5 lev 25x (IM 4.0%, MM 0.5%)
liquidation price = entry*(1-IM)/(1-MM) = 48,241
drop to liq = 3.52%  vs naive 1/leverage = 4.00%
check: loss 879 = wallet 1,000 - MM 121 ? True
   2x: liq 25,126  (49.75% away)
   5x: liq 40,201  (19.60% away)
  10x: liq 45,226  (9.55% away)
  25x: liq 48,241  (3.52% away)
  50x: liq 49,246  (1.51% away)
\end{lstlisting}
\begin{solution}
Equation \eqref{eq:liq} is checked numerically: loss $q(E-P_{\mathrm{liq}})=879$ equals wallet
$1{,}000$ minus maintenance value $121$ at the liquidation price. Note the check uses MM on the
\emph{liquidation} notional $qP_{\mathrm{liq}}$, not the entry---a one-line bug that shifts the
answer by bp and shows up in exercise~\ref{ex:liqex}.
\end{solution}
\end{example}

\begin{rulebox}[On venues (rules change)]
Perp mechanics vary: funding epoch (1h Hyperliquid vs 8h legacy venues), premium smoothing,
mark price (index vs impact-mid vs fair-price formulas), margin modes (cross vs isolated),
and whether liquidation is partial or full. Quote a venue rule only with a date and a re-check
plan; treat every number in this chapter's venue claims as \emph{[unverified]} until you check
the live docs.
\end{rulebox}

\section{Funding as data}\label{sec:funddata}

Persistent positive funding means longs are paying to stay---crowded bullish positioning and a
fragile squeeze candidate. Persistent negative funding is the mirror. Funding is also a yield
instrument: the mean rate annualises to a carry estimate. From our 500h snapshot:
mean $+0.113$\,bp/h $\Rightarrow$ about $+9.9\%$/yr for the short side---before fees, before
flips, and before the basis moves against you.

\begin{example}\label{ex:fundread}
\textbf{Reading a tape.} BTC rallies 6\% over a day while funding runs $+0.8$\,bp/h for
sixteen consecutive epochs (versus the $0.125$bp interest baseline). Open interest is up 12\%.
What is the tape saying, and what is the carry cost of staying long 1\,BTC at index 60{,}000?
\begin{solution}
Funding well above the interest floor says the premium term is engaged: the perp trades
persistently rich, i.e.\ buyers are crowding the levered long side and paying for the privilege.
Rising open interest alongside means new leveraged longs, not just spot-driven move---the
classic crowded-long squeeze setup. Cost: $60{,}000\times0.8\times10^{-4}\times24 \approx
\$115$/day per BTC held long. A trader who stays long is paying $\approx1.7\%$/month to hold
the position, which is the position's break-even drift, not a footnote.
\end{solution}
\end{example}

\begin{example}\label{ex:dash}
\textbf{The funding dashboard.} \texttt{funding\_hist.py} is the whole skill in 20 lines:
fetch (or load pinned) funding, then report mean rate, share of positive epochs, distance from
the interest floor, and an annualised carry. Its output is quoted in \S\ref{sec:funding}.
Build the same tool against any venue's public funding endpoint and you have a positioning
gauge that needs no price model at all.
\end{example}

\begin{checkpoint}[Perp math]
\begin{itemize}
\item Derive \eqref{eq:liq} for a \emph{short}: $P_{\mathrm{liq}} = E(1+\mathrm{IM})/(1+\mathrm{MM})$---sign flip, not a new formula.
\item In \texttt{data/hl\_funding\_btc.json}, find the longest run of $f_t<0$ and compute what a 10\,BTC short earned that day.
\item Why does the interest floor, not the premium, dominate mean funding in quiet regimes? (Answer: the clamp zeroes the premium term while the floor always pays.)
\end{itemize}
\end{checkpoint}

\begin{example}[What the window actually costs]\label{ex:carrycost}
\textbf{Strategy.} Take the pinned 500-hour BTC funding tape and
bill it like a broker would: multiply the summed rate by notional,
then ask how unevenly the cost arrives.

\textbf{Solution.} \texttt{carry\_cost.py} prints:
\begin{verbatim}
window: 500h of BTC funding on Hyperliquid
mean funding = 0.113 bp/h  (9.88% per yr if sustained)
mean premium = -0.965 bp/h
median funding = 0.125 bp/h  at-floor hours 85%
carry on $100k long over 500h = $563.83 = 56.4 bp
worst 8h funding stretch = 1.00 bp  best = -0.34 bp
\end{verbatim}
A $\$100$k long held the whole window pays $\$563.83$ ---
$56.4$bp, or an annualised $9.9\%$ if this rate sustained. Yet
85\% of hours sat at the $0.125$bp interest floor: the carry is a
\emph{fixed toll}, not compensation for a real premium (mean
premium $-0.97$bp --- the perp actually traded \emph{below} index
most of the window).

\textbf{Significance.} Funding is rarely ``the carry trade's
yield'' --- on this venue it is mostly the rate floor. The worst
8h stretch ($1.0$bp) is five times the mean: sizing a position
by average funding understates the crowded-window cost by the
same factor.
\end{example}

\section{Recap: perp math end-to-end}\label{sec:perprecap}

The chapter's toolkit is deliberately small. A perp position is fully described by
$(q, E, \mathrm{IM}, \mathrm{MM})$; its PnL is $q(P_t-E)$ minus accumulated
$qN_t f_t$; its survival is governed by \eqref{eq:liq}. Funding converts crowdedness into a
price series you can download and study. That is the entire engine---everything else a venue
shows you is presentation of these same numbers.

\begin{figure}[htbp]\centering
\begin{tikzpicture}[every node/.style={font=\scriptsize,align=center}]
  \node[draw=good,fill=goodbg,rounded corners=1pt,minimum width=3.4cm,minimum height=0.9cm] (spot)
    {long 1 BTC spot};
  \node[draw=amber,fill=amberbg,rounded corners=1pt,minimum width=3.4cm,minimum height=0.9cm,right=2.6cm of spot] (perp)
    {short 1 BTC perp};
  \node[draw=slate,fill=graybg,rounded corners=1pt,minimum width=3.4cm,minimum height=0.9cm,
        below=1.0cm of $(spot)!0.5!(perp)$] (res)
    {price moves cancel:\\ earns $f_t$ when $f_t>0$, pays when $f_t<0$};
  \draw[<->,ink,thick] (spot) -- node[above]{delta $=0$} (perp);
  \draw[->,ink] (spot) -- (res);
  \draw[->,ink] (perp) -- (res);
\end{tikzpicture}

\caption{The cash-and-carry structure: equal and opposite delta, funding as the return stream.}
\label{fig:carrystack}
\end{figure}

\begin{table}[htbp]\centering\small
\caption{Chapter toolbox.}
\begin{tabular}{@{}lll@{}}
\toprule
Concept & Formula/tool & Gotcha\\
\midrule
Funding payment & $qN_t f_t$ per epoch & sign convention varies by venue\\
Basis & $F/S-1$, annualise $\times365/D$ & perp basis has no expiry anchor\\
Carry & $\sum f_t$ on short-perp leg & flips sign; liq risk on short leg\\
Liquidation & \eqref{eq:liq} & MM on liquidation notional\\
Mark vs last & PnL on mark, trades on last & liquidations trigger off mark\\
\bottomrule
\end{tabular}
\end{table}

\section{Common mistakes}
\begin{itemize}
\item \textbf{``$1/\text{leverage}$ is my liquidation distance.''} Maintenance margin pulls it
  closer: $25\times$ at $\mathrm{MM}=0.5\%$ is $3.5\%$, not $4\%$.
\item \textbf{Forgetting funding in PnL.} A flat market is not free: at $+0.5$\,bp/h a long
  bleeds $\sim1.2\%$ of notional per day.
\item \textbf{Confusing mark and last price.} Liquidations, margin checks and displayed PnL
  use the mark; your fills are at last. In a wick, the difference is the whole story.
\item \textbf{Treating carry as riskless.} Short perp $+$ long spot still has liquidation
  risk on the short leg, funding-flip risk, and basis risk while entering/exiting.
\item \textbf{Reading funding as a signal after it prints.} Funding tells you who is
  \emph{already} crowded; it is a positioning measure, not a forecast.
\item \textbf{Annualising from one epoch.} $f_t$ is noisy at hourly scale; annualise the
  \emph{mean} over a window, not a single hot print.
\item \textbf{Annualising the floor.} $0.125$bp/h $\times
24\times365 = 10.95\%$/yr --- but that is the \emph{floor}, not a
signal: 85\% of the window printed exactly it. Quote the floor and
the premium separately.
\item \textbf{Budgeting carry by the mean.} The worst 8h stretch
cost $5\times$ the hourly mean (Example~\ref{ex:carrycost}). A
position sized to ``average funding'' is sized wrong precisely
when it matters most.
\end{itemize}

\section{Exercises}
\begin{exercise}[easy]
At index 80{,}000 you hold 0.4\,BTC long. Funding prints $+0.2$\,bp/h for 24 hours. Total paid?
\end{exercise}
\begin{solution}
$0.4\times80{,}000\times0.2\times10^{-4}\times24=\$15.36$.
\end{solution}

\begin{exercise}[easy]
Recompute $P_{\mathrm{liq}}$ for the entry-50{,}000 long at $10\times$ ($\mathrm{IM}=10\%$) with
$\mathrm{MM}=0.5\%$. Then for $5\times$.
\end{exercise}
\begin{solution}
$10\times$: $50{,}000\times0.9/0.995=45{,}226$ ($9.5\%$ away).
$5\times$: $50{,}000\times0.8/0.995=40{,}201$ ($19.6\%$ away)---matching the table.
\end{solution}

\begin{exercise}[easy]
A perp trades at index $+45$\,bp. Funding is currently $+0.12$\,bp/h. Is funding doing its job,
and what would you expect to print over the next hours?
\end{exercise}
\begin{solution}
A persistent $+45$bp premium should raise the premium term: expect $f_t$ to climb above the
interest floor toward the venue's clamp, paying shorts to lean against the rich perp until the
premium compresses.
\end{solution}

\begin{exercise}[easy]
Your venue quotes funding ``0.01\%'' per 8h. What is that in bp/h and in \%/yr at a flat rate?
\end{exercise}
\begin{solution}
$0.01\%=1$\,bp per 8h $=0.125$\,bp/h; flat: $1\,\mathrm{bp}\times3\times365/10^{4}=10.95\%$/yr.
\end{solution}

\begin{exercise}[easy]
Short 1\,BTC at entry 60{,}000, $20\times$ ($\mathrm{IM}=5\%$), $\mathrm{MM}=0.5\%$. Liquidation
price?
\end{exercise}
\begin{solution}
Shorts liquidate on rallies: $P_{\mathrm{liq}}=E(1+\mathrm{IM})/(1+\mathrm{MM})
=60{,}000\times1.05/1.005\approx62{,}687$ ($4.5\%$ above).
\end{solution}

\begin{exercise}
From \texttt{data/hl\_funding\_btc.json}: compute the maximum 24h window sum of $f_t$, in bp,
and the carry a 5\,BTC short would have earned in that window at notional \$50{,}000.
\end{exercise}
\begin{solution}
Rolling-sum the pinned series (max window $\approx$ sum of $0.125$bp/h plus premium bursts,
$\approx3.5$bp/day depending on hours counted); carry $=5\times50{,}000\times3.5\times10^{-4}
\approx\$87.5$. Your script's exact numbers are the answer.
\end{solution}

\begin{exercise}
Derive \eqref{eq:liq} for cross margin where wallet $w$ exceeds IM on the position. What
changes?
\end{exercise}
\begin{solution}
Replace IM by $w/(qE)$ in \eqref{eq:liq}: more collateral moves liquidation further away;
isolated is just cross with $w=$ IM$\times$notional.
\end{solution}

\begin{exercise}\label{ex:liqex}
The chapter warns MM applies to liquidation notional. Compute the liq price if you
(incorrectly) apply MM to entry notional: $P^*=E(1-\mathrm{IM}+\mathrm{MM})$. At $25\times$,
what error does that make in the distance-to-liq?
\end{exercise}
\begin{solution}
$P^*=50{,}000(1-0.04+0.005)=48{,}250$ vs correct $48{,}241$: the naive formula lands $9$\,USD
closer, a $0.02$\,bp... about $0.9$\,bp of price---small at entry scale, but it is the
difference between a checked margin engine and a moved one.
\end{solution}

\begin{exercise}
In \S\ref{sec:funding} we saw 97\% of hours had $f_t>0$ while mean premium was $-1.0$\,bp.
Reconcile.
\end{exercise}
\begin{solution}
The interest component pays longs-to-shorts regardless of premium sign; only when the
clamped premium term is sufficiently negative ($<-0.125$bp/h) does $f_t$ flip. Premium $-1$bp
on the impact-mid definition is near zero after the venue's smoothing.
\end{solution}

\begin{exercise}
A carry desk quotes ``30\% APR available on BTC funding.'' Sketch the position, list its three
independent failure modes, and state which one is capped by venue design.
\end{exercise}
\begin{solution}
Long spot $+$ short perp. Failure: (i) funding flips negative---uncapped, continuous risk;
(ii) liquidation of the short leg during basis widening---capped only by posting collateral,
i.e.\ the trade can be right and still forced out; (iii) basis moves against you at
entry/exit. None is capped by venue design except the rate clamp $f_{\max}$ which bounds the
\emph{speed} of your bleed, not the side.
\end{solution}

\begin{exercise}
Show that for a perp held $D$ days, funding's drag on a levered long's \emph{return on margin}
is $f_{\mathrm{day}}\times D\times\mathrm{lev}$ where $f_{\mathrm{day}}$ is the daily funding
rate as a fraction of notional.
\end{exercise}
\begin{solution}
Daily drag on notional is $f_{\mathrm{day}}$; margin is notional/lev, so drag on margin
$=f_{\mathrm{day}}\times$lev per day, times $D$ days.
\end{solution}

\begin{exercise}
Funding is a peer transfer. So why does the venue care about the sign? What breaks if $f_t$
sticks at $+5$\,bp/h for a week?
\end{exercise}
\begin{solution}
Extreme funding bleeds longs into margin calls, forcing deleveraging---the venue cares because
its risk engine inherits the position when margin fails. Sustained extreme rates are
themselves a solvency-stress signal.
\end{solution}

\begin{exercise}[hard]
Suppose impact-mid premium $p_t$ is AR(1) mean $\mu$, persistence $\phi$. Show that
$\mathrm{E}[f_t] \to \mathrm{clamp-free}$ value $I+\mu/(1-\phi)$ is wrong---derive the correct
mean and explain why the clamp matters even when rarely binding.
\end{exercise}
\begin{solution}
The MA form is $\tilde p_t=\mu/(1-\phi)+\text{noise}$; the clamp acts on the \emph{rate},
truncating its right tail, so $\mathrm{E}[f] < I+\mathrm{E}[\tilde p]$ whenever the upper tail
crosses $f_{\max}$---the bias exists even at low binding frequency because truncation removes
mass without replacing it.
\end{solution}

\begin{exercise}[hard]
Extend \texttt{funding\_sim.py}: position switched to \emph{short perp $+$ long spot}. Report
realised carry distribution over 200 seeds and the fraction that lose money. Discuss.
\end{exercise}
\begin{solution}
Modify the file to flip sign of funding receipts and add the spot leg; losers are seeds where
premium mean-reverts below $-I$ for sustained windows. The exercise is the point: implement it
and report your number (typical runs show a few percent of seeds net-negative at these
parameters).
\end{solution}

\begin{exercise}[hard]
Funding $f_t$ appears both as a transfer and as a forecasting variable. Show why
$\mathrm{E}[r_{t+1}\mid f_t] \neq 0$ does \emph{not} imply a free edge, using the crowding
interpretation.
\end{exercise}
\begin{solution}
Positive predictability (high funding predicting mean-reversion/squeezes) is exactly what
keeps the mechanism at its resting point: if a simple rule printed money, crowding into the
rule would move the premium and the funding equilibrium. The variable is a thermometer, and
trading the thermometer moves the temperature.
\end{solution}

\begin{exercise}[easy]
On the pinned window a $\$100$k long pays $\$563.83$ over 500h.
What would a $\$100$k \emph{short} have received, and why is the
answer not exactly $-\$563.83$?
\end{exercise}
\begin{solution}
The same magnitude --- funding is symmetric in sign. It is not
exactly $-\$563.83$ only because the rate floor makes some hours
structurally nonzero regardless of premium sign.
\end{solution}
\begin{exercise}[medium]
The mean premium was $-0.97$bp/h while mean funding was
$+0.113$bp/h. Reconcile the signs in one sentence.
\end{exercise}
\begin{solution}
Premium negative means perp $<$ index; the rate floor still
charges longs the default interest component, so sign(funding)
$\neq$ sign(premium) whenever the premium is small.
\end{solution}
\begin{exercise}[hard]
Your strategy shorts perps to ``collect funding''. Using the pinned
window, estimate P\&L per $\$100$k short including the premium drift
you'd also earn/pay, and state the assumption that breaks the
trade.
\end{exercise}
\begin{solution}
Collect $+\$563.83$ funding but pay the negative premium drift
($-0.97$bp/h $\times$ 500h $\times \$100$k $= -\$485$): net
$\approx +\$79$. The break: any window where premium reverts
positive flips the drift against you --- you are short the perp
basis, not long a coupon.
\end{solution}

\section*{Notes}
Funding is the invisible engine of crypto leverage: it prices crowdedness continuously and it
is the reason ``free'' leverage is never free. The margin engine and \eqref{eq:liq} are the
other half---most retail losses are liquidation mechanics, not bad calls. For live rules see
venue docs (they change); for the market context, the CFTC's 2025 perpetual-futures event is
the regulator's own explainer of the same machinery.

\section*{Further reading}
\begin{itemize}
\item Hyperliquid docs, ``Funding'' --- the venue's own formula, clamp and interest terms:

\url{https://hyperliquid.gitbook.io/hyperliquid-docs/trading/funding}
\item Wikipedia, ``Perpetual futures'' --- history (Shiller) and the mechanism family:

\url{https://en.wikipedia.org/wiki/Perpetual_futures}
\item J.\ Hull, \emph{Options, Futures, and Other Derivatives} --- the expiry-convergence
futures theory this chapter leans on:

\url{https://www.pearson.com/en-us/subject-catalog/p/futures-options-and-other-derivatives/P200000002816}
\item CFTC, AAC meeting on perpetual futures (2025) --- the regulator's explainer, incl. the
funding-rate session video:

\url{https://www.cftc.gov/PressRoom/Events/opaeventtechperp072925}
\item L.\ de Prado, lectures --- for carry and positioning read in a research frame:

\url{https://www.quantresearch.org/Lectures.htm}
\end{itemize}

\begin{rulebox}[Watch instead]
\begin{itemize}
\item ``Presentation on Perpetual Futures' Funding Rate Mechanism'' --- CFTC (14:31). The
mechanism from the regulator: no-expiry design, funding replacing convergence.
\url{https://www.youtube.com/watch?v=ZEPjpgVSavY} (accessed 2026-10-11)
\item ``How Smart Traders Use Funding Rates'' --- Trading Bureau (6:45). Funding as
positioning data; matches \S\ref{sec:funddata}.
\url{https://www.youtube.com/watch?v=t06z4rCmp4w} (accessed 2026-10-11)
\end{itemize}
\end{rulebox}
